\chapter{Conclusões e Trabalho Futuro}
\label{chap:conc-trab-futuro}
\section{Conclusões Principais}
\label{chap5:sec:conclusoes_principais}

A proliferação de dispositivos da \textit{Internet of Things} (IoT) em ambientes críticos, como as \textit{Smart Cities}, exige uma mudança de paradigma na cibersegurança. Este projeto propôs, desenvolveu e validou uma arquitetura de Prevenção de Intrusões (IPS) autónoma, deslocando a inteligência analítica da nuvem para a periferia da rede (\textit{Edge Computing}).

A principal conclusão extraída deste trabalho é que a viabilidade de um sistema de segurança preditivo operando num \textit{Edge Router} não depende exclusivamente da precisão estatística máxima, mas sim de um compromisso rigoroso entre a eficácia preditiva e a latência operacional. Ao avaliar empiricamente quatro algoritmos de \textit{Machine Learning} distintos, provou-se que modelos baseados em árvores de decisão profundas (como o \textit{Random Forest}) geram latências e flutuações de processamento incompatíveis com o encaminhamento de tráfego em tempo real. 

Em contrapartida, o modelo \textit{Multi-Layer Perceptron} (MLP) revelou-se a solução ótima para este cenário. A sua arquitetura matricial permitiu classificar anomalias complexas com um tempo de inferência ultrarrápido (próximo de 1 milissegundo por pacote) e com uma variância praticamente nula. 

Adicionalmente, a implementação de uma topologia de emulação realista (Mininet) demonstrou o sucesso da abordagem híbrida do sistema. A conjugação da inteligência do MLP com o bloqueio a nível de rede (através de uma \textit{Blacklist} dinâmica) provou ser capaz de mitigar injeções massivas de tráfego malicioso, como ataques de Negação de Serviço (DDoS) e \textit{Botnets} (Mirai). O sistema descartou eficazmente as ameaças nas primeiras camadas da rede, preservando o poder computacional do \textit{hardware} e garantindo o funcionamento ininterrupto dos dispositivos legítimos da cidade.

\section{Trabalho Futuro}
\label{chap5:sec:trabalho_futuro}

Apesar dos resultados promissores alcançados, a segurança em redes IoT é um campo em constante mutação. A arquitetura desenvolvida abre portas a diversas linhas de investigação futura para robustecer o sistema.

Em primeiro lugar, a abordagem de extração de \textit{features} baseada em fluxos estatísticos, embora fundamental para garantir a baixa latência e a privacidade dos dados, limita a deteção de anomalias ao nível da aplicação (Camada 7) ou ataques altamente furtivos. Trabalhos futuros poderão explorar a integração de Inspeção Profunda de Pacotes (DPI) de forma seletiva, em que o \textit{payload} só é analisado caso o modelo MLP detete um comportamento estatisticamente suspeito, criando assim um mecanismo de verificação em duas etapas.

Em segundo lugar, seria de elevado interesse integrar este nó \textit{Edge} com controladores de redes definidas por \textit{software} (\textit{Software-Defined Networking} - SDN). Num cenário de \textit{Smart City}, se um nó intercetar um atacante \textit{Zero-Day}, o controlador SDN poderia propagar a \textit{Blacklist} gerada a todos os restantes \textit{routers} da cidade instantaneamente, criando um sistema imunológico distribuído e automatizado.

Por fim, a expansão desta arquitetura poderia beneficiar da adoção de algoritmos de \textit{Federated Learning} (Aprendizagem Federada). Esta tecnologia permitiria que múltiplos servidores \textit{Edge} espalhados por diferentes cidades treinassem e melhorassem os seus modelos de Inteligência Artificial de forma colaborativa, partilhando apenas as atualizações matemáticas do modelo, sem nunca transferir o tráfego sensível dos cidadãos para um servidor central, elevando assim os padrões de privacidade e segurança da infraestrutura ciberfísica.